\documentclass[10pt,a4paper,twocolumn]{article}

% FONT AND ENCODING
\usepackage[utf8]{inputenc}
\usepackage[english]{babel}
\usepackage[T1]{fontenc}            % Important for correct hyphenation
\usepackage[scaled=0.85]{beramono}  % Good mono font for code
\usepackage{newtxtext,newtxmath}    % Replaces obsolete 'times' package
\usepackage{csquotes}               % Context sensitive quotes
\usepackage{amsmath}                % (Note: newtxmath loads amssymb automatically)
\usepackage{minted}

% GEOMETRY
\usepackage[a4paper, margin=2cm, columnsep=0.6cm]{geometry}

% TABLES AND FIGURES
\usepackage{graphicx}
\usepackage{subcaption}
\usepackage{tabularx}
\usepackage{multirow}
\usepackage{float}
\usepackage{placeins}

% BIBLIOGRAPHY
\usepackage[numbers,sort&compress]{natbib}
\bibliographystyle{IEEEtran}

% HYPERLINKS
\usepackage{hyperref}
\usepackage[capitalize,noabbrev]{cleveref}

% CUSTOM
\setlength{\parindent}{1em}         % Academic papers usually indent paragraphs
\setlength{\parskip}{0pt}           % And have no skip between them
\title{%\begin{center}
%\includegraphics[width=2.56in]{imgs/logo.png}\\
    %{{\Large{\textsc{Alma Mater Studiorum $\cdot$ Università di
    %Bologna}}}} \rule[0.1cm]{14.7cm}{0.1mm}
    %\rule[0.5cm]{14.7cm}{0.6mm}
    %\vspace{5mm}
%\end{center}
    \LARGE{\bf Predictive Deception: \\
    \vspace{3.5mm}
    LLM-based Command Anticipation\\
    \vspace{3.5mm}
    in SSH Honeypots\\
    \vspace{3.5mm}}}
\author{Leonardo Barone\\
Leonardo Ciacco\\
Enrico Giannini}
\date{} %nasconde la data, si può inserire \date{December 17, 2025}
%aggiunta per evitare che il link sbrodoli
\def\UrlBreaks{\do\/\do-}
\begin{document}
\setcounter{page}{1} 
\pagenumbering{arabic}
%%%%%%%%%%%%%%%%%%%%%%%%%%%%%%%%%%%%%%
% TITLE PAGE
%%%%%%%%%%%%%%%%%%%%%%%%%%%%%%%%%%%%%%


\twocolumn[
    \maketitle
    \begin{@twocolumnfalse}
        \renewcommand{\abstractname}{Abstract}
        \phantomsection
        \begin{abstract}
            Traditional SSH honeypots operate on a reactive paradigm, responding to attacker inputs only after execution, which limits their ability to shape the attack surface dynamically. This research introduces \textit{Predictive Deception}, a framework that transforms honeypots from passive sensors into intelligent, proactive adversaries. By leveraging Large Language Models (LLMs) to analyze command sequences in real-time, the proposed system anticipates attacker's intent before specific actions are taken. Historical attack data are processed using a sliding window approach, and subsequently used for fine-tuning an LLM with the purpose of being precise in identifying the most probable purposes of the attack, hence, the final command.
            These predictions drive an automated orchestrator that proactively injects deceptive artifacts, taking advantage of \textit{Canarytokens}, into the environment immediately prior to the attacker’s request. This approach aims to minimize detection latency and maximize engagement time, effectively keeping the system ahead of the adversary while making it think the system is legit and not a decoy. This study evaluates whether such probabilistic anticipation provides higher-fidelity threat intelligence compared to standard reactive deception.
            %The research question is whether an LLM  can predict an attacker's next command with sufficient accuracy to enable proactive placement of deceptive artifacts in SSH honeypots, and and if this does increase detection value compared to reactive-only deception.\\
            \vspace{1cm}
        \end{abstract}
    \end{@twocolumnfalse}
]
\setcounter{secnumdepth}{-1}
% MAIN
\raggedbottom
\pagestyle{plain}
% TODO: riguardare struttura e contenuti + modificare commenti, inoltre sistemare titoli

% INTRO
\section{Introduction}
This section formalizes the proposed \textit{Predictive Deception} approach and the experimental protocol used to evaluate it. The core idea is to anticipate an attacker's actions in the near future from SSH session histories and to proactively deploy deception artifacts (Canarytokens) before those actions are executed.

\subsection{Goals and hypotheses}
We focus on two main questions:
\begin{enumerate}
    \item \textbf{(Prediction)} can an LLM infer an attacker's near-future actions from an observed command prefix in a live SSH session?
    \item \textbf{(Deception efficacy)} does proactive placement of deception artifacts improve detection and engagement compared to reactive-only deception?
\end{enumerate}
We test the hypotheses that predictive deception decreases time-to-detection (first canary trigger) and increases deception interaction rate compared to reactive-only strategies.

\subsection{Contributions}
This work makes the following contributions:
\begin{itemize}
    \item a predictor-executor architecture that anticipates near-future attacker actions from live SSH command prefixes and proactively deploys deception artifacts;
    \item a horizon-based prediction formulation made for SSH command streams, using a sliding window and top-$K$ anticipated high-value actions rather than a simpler, but too slow and possibly inaccurate, exact-string next-command prediction;
    \item an evaluation protocol that compares predictive deception against a reactive baseline on two public datasets, reporting both prediction quality and deception outcomes.
\end{itemize}

\subsection{Report Organization}
The report is organized as follows: \hyperref[sec:state-of-the-art]{Background and State of the Art} summarizes related concepts and tools; \hyperref[sec:datasets]{Datasets} describes the sources and preprocessing goals; \hyperref[sec:methodology]{Methodology} formalizes the prediction task and system design; \hyperref[sec:model-impl-training]{Model Implementation and Training} details dataset preparation, integration, enrichment and training; \hyperref[sec:experiments-results]{Experiments and Results} reports baselines, metrics and outcomes; \hyperref[sec:threats-validity]{Threats to Validity and Reproducibility} discusses limitations.

\section{Background and State of the Art}
\label{sec:state-of-the-art}

\subsection{Reactive SSH Honeypots}
A \textbf{honeypot} is a system or software designed to simulate a real environment, attracting attackers while logging and analyzing their activities.\\
There happens to be a wide variety of honeypots, distinguished by field of application, size or interaction level. The latter metric varies from \textbf{low} (simple protocols,
undetailed informations, but easy to set up) to \textbf{medium} and \textbf{high} levels of interaction (widening of services provided along with greater costs) \cite{Honeypots101}.\\

\textbf{Cowrie} is a medium to high (depending on the configuration) interaction SSH and Telnet honeypot designed to log brute force attacks and the shell interaction performed by the attacker.

On default, it sets up a fake fylesystem through a python infrastructure that provides a facade that resebles a Debian 5.0 installation. 
This tool will be further used to build a \textbf{predictive behaviour} on top of a static finite-state-machine one \cite{CowrieDocumentation}.\\

\subsection{Canarytokens}
Canarytokens are deployed inside the file system and, just like sensors, they notify every bit of movement around them. They are meant to be lightweight tools that disguise themselves as interesting and appealing information, such as:
\begin{itemize}
    \item \textbf{AWS API Keys}: this set of fake cloud infrastructure keys will send a notification as soon as they are used;
    \item \textbf{Kubeconfig}: a .yml file that should contain all the information needed to authenticate and control a Kubernetes cluster, but instead returns permission errors and sends an alert to whoever placed the trap;
    \item \textbf{Adobe PDF}: this is a PDF file that forces the reader to do a DNS lookup to the Canarytoken server when opened, notifying the intrusion;
    \item \textbf{WireGuard token}: it fakes a WireGuard VPN endpoint so that, when the machine is compromised and the attacker tries to enumerate the VPN configurations and connect to them, it sends an alert;
    \item \textbf{HTTP Canarytoken}: this canarytoken generates an URL that triggers an alert when visited, it will be used in this project to replace malware files with non-malicious scripts and executables, to trick the attacker into revealing its presence.
\end{itemize}
\cite{CanarytokensGuide}

\subsection{LLM fine-tuning}
The efficacy of Large Language Models in offensive cybersecurity, was demonstrated largely, for example with tools like \textbf{PentestGPT} utilizing LLM-driven reasoning to automate complex penetration testing tasks \cite{PentestGPT}. To counter such sophisticated, adaptive threats, defensive systems must evolve beyond static interactions. \textbf{Fine-tuning} provides the necessary bridge, applying transfer learning to adapt a pre-trained LLM into a specialized defensive component capable of simulating realistic Unix environments, with a better chance of deceiving the attacker.\\ % TODO: modify this section

However, generic models often suffer from \textbf{hallucinations}, generating syntactically correct but functionally incoherent responses, which can quickly reveal the deception to an attacker. Fine-tuning addresses this by training the model on targeted datasets of real system logs and command-line interactions. This process is generally approached through two methodologies:
\begin{itemize}
    \item \textbf{Full Fine-Tuning}: updates all the parameters of the model. While it offers the highest adaptability, it is computationally expensive and resource-intensive;
    \item \textbf{Parameter-Efficient Fine-Tuning (PEFT)}: techniques, such as \textbf{LoRA} (Low-Rank Adaptation), preserve the original model weights and add small, optimizable low-rank matrices into each Transformer layer. This lets us fine-tune the model without updating most of its parameters by just updating the added matrices. As a result, training becomes much faster, uses far fewer resources, and is easier to deploy, even on lightweight systems like Cowrie. \cite{LLMFineTuning}.
    \item \textbf{Supervised finetuning}, for example for instruction finetuning to teach a model to follow instructions. Requires a dedicated dataset of prompts/responses.
\end{itemize}
The ultimate goal is to transition from a static system to a \textbf{dynamic response engine}. Exploiting the patterns in past attacks, the honeypot can proactively inject deceptive artifacts, prolonging engagement, improving data collection and increasing its success rate in making the attacker interact with canary files. 

% TODO: controllare se espandere la sezione o va bene così
\section{Datasets}
\label{sec:datasets}
We use two publicly available SSH honeypot datasets \cite{CyberLab_JSON_DataSet,CowrieHoneypot_DataSet}. One dataset is used as primary training data (large-scale Cowrie logs), while the other is used as an external evaluation set to test cross-domain generalization (smaller, independently collected sessions). Both datasets are preprocessed to extract ordered command sequences per session, along with timestamps for temporal analysis (see \hyperref[sec:dataset-prep]{dataset preparation section}).

The research is carried out on the basis of two data sets \cite{CyberLab_JSON_DataSet} \cite{CowrieHoneypot_DataSet}. Those work as a training base \footnote{the latter}, and as a valid counterpart \footnote{the former, considering it has a notably smaller set of attacks}, in order to measure the relative performance of predictive deception models to the reactive ones.\\

\section{Methodology}
\label{sec:methodology}

\subsection{Strategy for Command Prediction}
An attacker session is modeled as an ordered command sequence $c_{1:n} = (c_1,\dots,c_n)$. At time step $k$, the system observes a history prefix $c_{1:k}$ and produces a prediction intended to support proactive trap placement.

Rather than predicting the single next command string (which is almost impossible to get right due to random filenames, IPs, and URLs), try to predict the final intent of the attacker. Given a fixed window size $m$, the model outputs a small set of anticipated \textit{high-value actions} expected to occur within the next $m$ steps, for example, download and execution attempts, credential discovery, enumeration of configuration files, or access to sensitive paths. In practice, the predictor returns the top-$K$ predicted actions (we use $K=3$) to balance coverage and latency.

\subsection{System Overview: Predictor--Executor Architecture}
%overview concettuale per poi introdurre il prossimo capitolo
The aim is to predict the attacker's next commands and to place the deceptive artifacts accordingly. In order to do so, the design is based on a \textbf{predictor - executor} model (Fig.\ref{fig:arch}) where the LLM Predictor anticipates the purpose of the attack, classifies it and gives it a label, subsequently the Executor properly places the trap files:

\begin{figure}[H]
    \centering
    \includegraphics[width=\linewidth]{imgs/Diagramma_logico_architettura.drawio.png}
    \caption{Logic architecture}
    \label{fig:arch}
\end{figure}

\subsection{Design Constraints and Online Inference Policy}
SSH honeypots are exposed to heterogeneous interaction patterns:
\begin{itemize}
    \item \textbf{Bot-driven sessions} (dominant in volume): short-lived, high command throughput, often fully automated.
    \item \textbf{Human-driven sessions}: longer sessions with larger inter-command delays and more exploratory behavior.
\end{itemize}
This motivates two design choices.
\textbf{(i) Sliding prediction horizon.} The predictor operates over a bounded horizon $m$ to limit computation and to focus on actions that can be exploited by deception (e.g., placing bait files before the attacker reaches them).
\textbf{(ii) Batched inference.} Instead of querying the model after every command, predictions are performed in batches (e.g., every $b$ commands or when a ``suspiciousness'' threshold is reached), reducing overhead.

Because model inference is slower than scripted attackers, we additionally introduce a small server-side delay to keep interactions realistic while giving the executor time to place artifacts. In our implementation, the added latency is sampled from a bounded distribution (e.g., $0.7\,\mathrm{s}$ plus a small random jitter). We report the exact distribution and its impact on measured engagement in \hyperref[sec:threats-validity]{Threats to Validity and Reproducibility}.

In practice, even the fastest LLMs cannot compete in terms of speed with scripted attackers that can issue dozens of commands quickly. For this reason, the predictor is queried as a one-shot call for each batch of observations, producing a small set of plausible branches the attacker could take; this keeps latency bounded while still enabling proactive template placement.

% TODO: add that we introduce 0.7 + [0..0.5] seconds of delay to manage this issue. rework this section



\section{Model implementation and Training}
\label{sec:model-impl-training}
\subsection{Dataset Preparation}
\label{sec:dataset-prep}
It was necessary to prepare scripts that initially parsed the CSV dataset \cite{CowrieHoneypot_DataSet} to a JSON format, and then to narrow down all the useless parameters, looking ahead to the training phase.\\
As presented by Fig. \ref{fig:dataprep} al data are regrouped into ordered sessions. 
\begin{figure}[H]
    \centering
    \includegraphics[width=0.5\linewidth]{imgs/diagramma dati.drawio.png}
    \caption{Dataset preparation course}
    \label{fig:dataprep}
\end{figure}
Sessions, made by N commands, will be then divided into N-1 subsessions with given future inputs, in order to highlight the association between the M future commands and the N-M history of inputs.\\

\begin{verbatim}
<considering 
session = {cmd_1, ..., cmd_m, ..., cmd_n}>
{
    "hystory": "cmd_1 ; cmd_2 ; ...; cmd_n-m ;" 
    "future": "cmd_n-m+1; ...; cmd_n"
}
\end{verbatim}
Regarding the translation from the "future" and the expected output of a prevision, two choices have been made: the adoption of a Sliding Window, and the filtering of the less-interesting commands predicted.
The idea of "Sliding Window" refers to  a maximum number of command seen in the near future. 
These two combined bring the focus of the expected results, not on the production of an exact command sequence, but rather on a group of high-relevance actions that the attacker does (such as looking for a specific file). The concept will not only be a pattern recognition but also a target prediction that the system will be able to exploit by placing the correct decoy. 
The result will be like asking: "which are the probable endpoint commands the attacker will bring in the next <sliding\_window> steps?".\\

In the end we have obtained all the dataset in a ready-to-train form such as:\\
\newpage
\begin{verbatim}
{
"instruction": "<detailed prompt>",

"input": "#!/bin/sh ; PATH=$PATH:/usr/local/sbin:
/usr/local/bin:/usr/sbin:/usr/bin:/sbin:/bin",

"output": "wget http://23.235.171.196:88/112 ;
curl -O http://23.235.171.196:88/112 ; 
chmod +x 112 ; ./112 ; rm -rf 123.sh"
}
\end{verbatim}
\subsection{System Integration: From Prediction to Deception} %parte di spiegazione appunto del passaggio dal concetto di predizione all'utilizzo dl templating
As it was said before, minimizing the LLM calls is necessary in order to maintain the \textbf{lowest latency} achievable.
Consequently, it's clear that a response that matches prediction with deception is crucial to transfer to an hypothetic orchestrator script every information it needs to put into place the deceptive strategy.
There could be various implementations, two of which are:
\begin{itemize}
    \item mapping every prediction to a dedicated deception pattern;
    \item creatively create a deception in response to the prediction through an LLM.
\end{itemize}
Neither of them seemed viable, so it has been decided to adopt an \textbf{hybrid solution} that could combine the \textbf{creative matching} that an LLM can provide and a more efficient pre-constructed and ready-to-use \textbf{template pool}. The key concept is to integrate the prediction and deception, associating the input command history with the output predicted effective strategy.\\
In  order to achieve this, it's important to:
\begin{enumerate}
    \item rework the training dataset replacing each predicted command with a JSON entry that has enough informations to allow a non-AI based Python script to fully operate the deception,
    \item prepare a template pool to which the LLM responses will be depending, that covers most of the attack scenarios and can provide a variety of plausible traps, those will be used by the Python script to efficently populate the file system. 
    \end{enumerate} 
\subsubsection{Template Arrangement}
Before any kind of action on the dataset it's best to define the templates, ready to be hydrated with the Canarytokens' parameters, which are strictly bound to the subset of Canarytoken types used in this project. The main strategies will then be:
\begin{itemize}
    \item \textbf{AWS API Keys}: a reactive trap, as said before, sends an alert when the fake AWS keys are used by the attacker;
    \item \textbf{Kubeconfig}: similar to an AWS trap, but related to a fake set of access information to a Kubernets cluster;
    \item \textbf{Adobe PDF}: various kinds of PDF files that may seem  interesting to an attacker, using a PDF Canarytoken to alert the attacker's presence; 
    \item \textbf{WireGuard token}:as the previous ones it's simply a Wireguard Canarytoken used in different scenarios to seem as realistic as possible;
    \item \textbf{Fake executable files}: .sh files that contain calls to dedicated HTTP Canarytokens among some other commands;
    \item \textbf{Malware replacement}: a file replacement strategy that prevents a malware installation and acts a substitution with a simple executable that launches an HTTP request to an already set up HTTP Canarytoken, tricking the attacker into running its own alert;
    \item \textbf{Simple text monitoring}: here the use of an HTTP Canarytoken is background, through a watcher script and a lookup table, the first monitors any occurrence of any simple file name from the table, inside the Cowrie logs, and alerts for every access on any of them. 
\end{itemize}
\subsection{LLM Based Data Enrichment}
The task the more expert LLM (teacher) had to complete is to read the command input history, read the future commands, and extract the three most probable ones form there. Then the predicted commands must be paired with a specific deceptive strategy from the pool inside of a JSON containing objects in this form:\\
\begin{verbatim}
    {
        "predicted_cmd": "<cmd>",
        "trap_path": "/path/to/trap",
        "trap_template": "<TRAP_TAG>"
    }
\end{verbatim}
It is highly important to redact a prompt able to get an \textbf{almost deterministic behaviour }from an LLM while still taking advantage of its capability of interpreting patterns.
In this project this is accomplished through a pair made by: an \textbf{accurate} and \textbf{detailed} system prompt, and a more \textbf{general} but \textbf{constraint focused} user task instruction file.\\

Associating such complex instructions to an already rich dataset could result in a discrete amount of computational power required. This is why the solution adopted makes use of Google VertexAI's batch elaborating system ($BatchPredictionJob$) \cite{google_vertex_ai_prediction_classes}, which provides an efficient tool specifically designed to take on large workloads through API calls. 

\subsection{Training}
To verify if this kind of approach to predictive deception is effective the chosen model are both from the Gemma 3 family: the 4 and 12 Billion parameters ones, pretrained version, both with a 4 bit quantisation. This allows to directly control if an expected more refined model can actually perform better in a \textbf{single task} fine tuning and assignment.\\
The single task matter can be counterintuitive: a model intended to be able to solve the largest variety of problems possible is now enrolled into an almost static behaviour, where the predictions it gives are restricted to a limited field. This, however, is bound to the already stated latency constraint, it can't be traded off in order to maintain resemblance to a real vulnerable system. In order to shape the model into a single task machine the two model are, as previously suggested, only the pretrained version, instead of the instruction tuned version, already accustomed to a general purpose chatty behaviour.\\
To maximize the outcomes the specific models were picked from the Unsloth Hugging Face model collection, since the 4 bit models have specific optimizations that enhance their capabilities. \cite{gemma_2025_report} \cite{gemma3_train_blog} \cite{unsloth_gemma3_4b_pt_unsloth_bnb_4bit}
\cite{unsloth_gemma3_12b_pt_unsloth_bnb_4bit}\\

\subsubsection{Data Formatting}
Even though these models are pretrained, using a gemma  tokenizer, already optimized to recognize the standard Gemma chat template patterns is the suggested approach. \cite{gemma_chat_template} 
The enriched data are subsequently integrated inside a couple of requests and responses, in order to correctly fine tune both gemma models as good as possible, the processi is then carried out thorugh the Unsloth framework.

\section{Cowrie Setup}

\subsection{API Server Configuration}

\subsection{Cowrie rework}

% TODO: modificare questa sezione e riguardarla
\section{Experiments and Results}
\label{sec:experiments-results}

\subsection{Baselines}
To evaluate the contribution of predictive deception, we compare against a \textbf{Reactive Baseline}.
This baseline consists of a standard Cowrie installation enhanced with specific reactive hooks on common download commands (such as \texttt{wget} and \texttt{curl}). In this configuration, deceptive artifacts (e.g., replacement executables) are deployed or triggered only after the command is issued, representing a traditional low-latency but purely reactive defense mechanism.

\subsection{Evaluation Metrics}
We report both \textbf{prediction quality} and \textbf{deception outcomes}.
\begin{itemize}
    \item \textbf{Prediction:} top-$K$ accuracy over action labels within the horizon; precision/recall/F1 for action classes; and (when evaluating raw commands) normalized matching after removing volatile tokens (IPs, URLs, random filenames).
    \item \textbf{Deception:} canary trigger rate per session; time-to-detection (first canary event from session start and from first suspicious command); and engagement proxies such as session duration and number of commands before disconnect.
\end{itemize}

% TODO: add experiments results here

\section{Threats to Validity and\\ Reproducibility}
\label{sec:threats-validity}
Potential threats include dataset shift between training and evaluation sources, label noise introduced by heuristic sessionization, and the possibility that artificial latency inflates engagement metrics. We mitigate these by evaluating on an external dataset, reporting normalization rules, and always comparing predictive and reactive strategies under the same latency policy. We also report model version, fine-tuning method, main hyperparameters, and random seeds to enable reproducibility.

\section{Conclusions and Future Work}


% BIBLIOGRAPHY
\bibliography{refs}

\end{document}